% Example file demonstrating how to use glossary terms
% Include this in your chapters to see how the usage works

% First use of a term should use \gls{} to reference it
A fundamental concept in quantum computing is the \gls{qubit}, which can exist in a 
state of \gls{superposition}. Unlike classical bits that can only represent either 0 or 1, 
qubits can represent both simultaneously.

% For subsequent uses, continue using \gls{} to maintain consistent behavior
When multiple \gls{qubit}s interact, they can become \gls{entanglement|entangled}, creating 
quantum correlations that have no classical counterpart. However, these quantum states 
are fragile and subject to \gls{decoherence} from environmental interactions.

% For plural forms, use \glspl{}
Quantum computers process information using \glspl{quantum-gate}, which manipulate qubits 
to perform calculations. One of the promising physical implementations uses \gls{trapped-ions}.

% To capitalize the first letter, use \Gls{} or \Glspl{} for plurals
\Gls{rsa} encryption is widely used in modern cryptography, but is vulnerable to attacks 
using \gls{shors-algorithm}. Similarly, \gls{ecc} provides efficient security but is also 
threatened by quantum computing.

% For acronyms that should be expanded on first use
\Gls{qkd} can provide theoretically secure key distribution, unlike classical key exchange 
methods that are vulnerable to quantum attacks. \Gls{bqp} represents problems that quantum 
computers can solve efficiently.

% Current efforts like the \gls{nist-pqc} aim to develop and standardize \gls{post-quantum-cryptography}, 
including promising approaches like \gls{lattice-based-cryptography} and \gls{hash-based-cryptography}.

% Note: When compiling, you need to:
% 1. Run pdflatex on your document
% 2. Run makeglossaries on your document
% 3. Run pdflatex again
% 
% Example:
% pdflatex main.tex
% makeglossaries main
% pdflatex main.tex 